% ============================================================
%  Sistema de Equações — Apostila Premium | PratiqueJá
%  Engine: XeLaTeX
% ============================================================
\documentclass[12pt]{article}
\usepackage{../../pratiqueja}

\newcommand{\hlm}[1]{{\color{darkblue}#1}}

\setsubject{Sistema de Equações}

\begin{document}
\pagenumbering{arabic}
\setstretch{1.2}
\color{bodytext}

\heroheader{Sistema de Equações}%
           {Substituição, adição, comparação e tipos de solução}%
           {Matemática · Avançado}

\setlength{\columnsep}{18pt}
\setlength{\columnseprule}{0.5pt}
\def\columnseprulecolor{\color{exborder}}

\begin{multicols}{2}

\TW{Def}{darkblue}{Definição}{O que é um sistema de equações?}

\regra{
Um \hl{sistema de equações} é um conjunto de duas ou mais equações que compartilham as mesmas variáveis. O objetivo é encontrar os valores que satisfazem \textbf{todas} as equações simultaneamente.

Focaremos nos \hl{sistemas lineares} — formados por equações do 1º grau. Cada equação representa uma reta, e a solução é o ponto de interseção entre elas.
}

\exemp{
\[
  \begin{cases}
    2x + 3y = 6 \\
    x - y = 2
  \end{cases}
\]
}

\TW{1}{badgepurple}{Tipos de Solução}{Consistente, dependente ou inconsistente}

\regra{
\begin{itemize}
  \item \textbf{Solução única} — sistema \hl{consistente e independente}. As retas se cruzam em um único ponto.
  \item \textbf{Infinitas soluções} — sistema \hl{consistente e dependente}. As equações representam a mesma reta (são múltiplos uma da outra).
  \item \textbf{Nenhuma solução} — sistema \hl{inconsistente}. Retas paralelas — nunca se cruzam.
\end{itemize}
}

\TW{2}{darkblue}{Método da Substituição}{Isolar uma variável e substituir na outra equação}

\regra{
Isola-se uma das variáveis em uma das equações e substitui-se na outra. Recomendado quando uma variável tem coeficiente 1 ou é fácil de isolar.
}

\exemp{
\[
  \begin{cases}
    x + y = 5 \quad \text{(1)} \\
    2x - y = 4 \quad \text{(2)}
  \end{cases}
\]

\textbf{Passo 1} — Isolar $x$ em (1): $x = 5 - y$

\textbf{Passo 2} — Substituir em (2):

$2(5-y) - y = 4$

$10 - 2y - y = 4$

$-3y = -6 \;\Rightarrow\; y = 2$

\textbf{Passo 3} — $x = 5 - 2 = 3$

$\hlm{x = 3 \;\;\text{e}\;\; y = 2}$
}

\TW{3}{darkblue}{Método da Adição}{Eliminar uma variável somando as equações}

\regra{
Multiplica-se uma ou ambas as equações por constantes de modo que, ao somá-las, uma das incógnitas seja eliminada.
}

\exemp{
\[
  \begin{cases}
    2x + y = 5 \quad \text{(1)} \\
    x - y = 1 \quad \text{(2)}
  \end{cases}
\]

Os coef. de $y$ são $+1$ e $-1$ — já opostos:

$(2x+y) + (x-y) = 5+1$

$3x = 6 \;\Rightarrow\; x = 2$

Substituindo em (2): $2 - y = 1 \;\Rightarrow\; y = 1$

$\hlm{x = 2 \;\;\text{e}\;\; y = 1}$
}

\obs{Quando os coeficientes não são opostos, multiplique uma ou ambas as equações por uma constante para torná-los opostos antes de somar.}

\TW{4}{darkblue}{Método da Comparação}{Isolar a mesma variável em ambas e igualar}

\regra{
Isola-se a \hl{mesma variável} nas duas equações e igualam-se as expressões obtidas, reduzindo o sistema a uma única equação com uma incógnita.
}

\exemp{
\[
  \begin{cases}
    7x + y = 140 \quad \text{(1)} \\
    -8x + y = -145 \quad \text{(2)}
  \end{cases}
\]

(1) $y = 140 - 7x$

(2) $y = -145 + 8x$

Igualando: $140 - 7x = -145 + 8x$

$285 = 15x \;\Rightarrow\; x = 19$

$y = 140 - 7 \cdot 19 = 140 - 133 = 7$

$\hlm{x = 19 \;\;\text{e}\;\; y = 7}$
}

\TW{Prat}{badgepurple}{Qual Método Usar?}{Guia rápido de escolha}

\regra{
\begin{itemize}
  \item Uma variável tem coef. 1 → \textbf{Substituição}
  \item Coeficientes de uma variável são opostos (ou facilmente tornam-se opostos) → \textbf{Adição}
  \item Fácil isolar a mesma variável nas duas equações → \textbf{Comparação}
\end{itemize}
Todos os métodos levam ao mesmo resultado — escolha o mais rápido para cada sistema.
}

\TW{Prob}{badgepurple}{Problemas Contextualizados}{Sistemas aplicados em situações reais}

\exemp{
\textbf{Prob. 1 — Idades:}

A soma das idades de Ana e Bruno é 30 anos. Bruno é 6 anos mais velho.

$\begin{cases} a + b = 30 \\ b = a + 6 \end{cases}$

Substituindo: $a + (a+6) = 30 \Rightarrow 2a = 24$

$\hlm{a = 12 \text{ anos}}$ \quad $\hlm{b = 18 \text{ anos}}$

\textbf{Prob. 2 — Ingresso:}

50 ingressos vendidos: inteira (R\$\,10) + meia (R\$\,5) = R\$\,380.

$\begin{cases} x + y = 50 \\ 10x + 5y = 380 \end{cases}$

Adição: mult. (1) por $-5$: $-5x - 5y = -250$

$(10x+5y) + (-5x-5y) = 380 - 250$

$5x = 130 \Rightarrow \hlm{x = 26}$ ingressos inteiros

$\hlm{y = 24}$ meias-entradas

\textbf{Prob. 3 — Velocidade:}

Trem A a 80 km/h e trem B a 100 km/h partem em sentido contrário a 360 km de distância.

$\begin{cases} 80t + 100t = 360 \end{cases}$

$180t = 360 \Rightarrow \hlm{t = 2 \text{ horas}}$
}

\end{multicols}

\end{document}
